\item \textbf{Método de desplazamiento para medir distancia focal.}
En un cuarto oscuro, se coloca una vela encendida a una distancia fija de $1.50\text{ m}$ de una pared blanca. Entre la vela y la pared se introduce una lente convergente en una posición tal que se proyecta una imagen nítida magnificada e invertida en la pared. La distancia al objeto en esta posición es $p_1$. Al desplazar la lente $90.0\text{ cm}$ hacia la pared, se forma una segunda imagen nítida en ella.
\begin{enumerate}
    \item Escriba la ecuación de la lente para la primera posición relacionando $f$ y $p_1$ únicamente.
    \item Escriba la ecuación para la segunda posición de la lente relacionando $f$ y $p_1$.
    \item Resuelva simultáneamente para encontrar la distancia objeto inicial $p_1$.
    \item Calcule la distancia focal $f$ de la lente convergente.
\end{enumerate}